% =====================================================================
%  11. Conclusion et perspectives
% =====================================================================
\section{Conclusion et perspectives}
\label{sec:conclusion}

Ce projet a traité la prédiction de la qualité de soudures d'acier comme un
ensemble de \textbf{régressions supervisées} sur des propriétés mécaniques
mesurées. L'analyse exploratoire a révélé une base compilée à partir de séries
expérimentales, avec des absences dépendant du procédé, des valeurs censurées,
des unités douteuses et des protocoles Charpy hétérogènes. Une préparation
entièrement traçable, un prétraitement appris dans chaque pli et une validation
croisée imbriquée selon deux scénarios ont permis une comparaison rigoureuse de
cinq familles supervisées et d'une méthode semi-supervisée.

Les résultats tiennent en trois phrases. \textbf{(1)} Avec 16 entrées et une
imputation médiane, la forêt aléatoire et la SVR RBF prédisent bien les soudures
de séries connues ($R^2$ de 0,66 à 0,84). \textbf{(2)} Le transfert à des séries
absentes reste difficile ($R^2$ de 0,08 à 0,34), principalement à cause de
décalages de niveau entre études que les entrées disponibles n'expliquent pas.
\textbf{(3)} COREG n'améliore pas la référence supervisée, car les entrées sans
label n'apportent pas l'information manquante.

\paragraph{Perspectives.}
\begin{itemize}
  \item \textbf{Validation indépendante} : évaluer la procédure retenue sur de nouvelles séries jamais utilisées pour la concevoir.
  \item \textbf{Enrichir les entrées} : réintégrer Nb, V, Ti, B ou Cu après vérification des unités par source, et tester leur apport par validation groupée.
  \item \textbf{Harmoniser Charpy} : poursuivre l'audit des 21 séries non résolues, séparer courbes $E(T)$ et températures de transition, puis construire une tâche homogène.
  \item \textbf{Quantifier l'incertitude} : intervalles de prédiction (forêts quantiles, prédiction conforme) et détection des entrées hors domaine.
  \item \textbf{Modèles hiérarchiques} : modéliser explicitement un effet aléatoire de série pour séparer variabilité intra- et inter-études.
  \item \textbf{Semi-supervisé} : tester COREG et d'autres méthodes dans un régime à labels rares (25 ou 50\,\%), où leur intérêt théorique est le plus grand.
\end{itemize}
